% ==============================================================================
% styles/code.tex
% ==============================================================================
\usepackage{minted}
\usepackage{etoolbox} % اضافه شده برای مدیریت هوشمند و امن محیط‌ها

\definecolor{codeBG}{HTML}{1E1E1E}
\definecolor{macRed}{HTML}{FF5F56}
\definecolor{macYellow}{HTML}{FFBD2E}
\definecolor{macGreen}{HTML}{27C93F}

\tcbuselibrary{minted, skins, breakable}

% تنظیم استایل شماره خطوط (انگلیسی و خاکستری)
\makeatletter
\renewcommand{\theFancyVerbLine}{%
  \sffamily\textcolor{gray}{\scriptsize\lr{\arabic{FancyVerbLine}}}%
}
\makeatother

% تنظیمات مشترک ظاهر مک
\tcbset{
    macbox/.style={
        enhanced, breakable,
        colback=codeBG,
        colframe=codeBG,
        colupper=white, % سفید کردن کاراکترهای عادی روی پس‌زمینه تیره
        arc=8pt,
        boxrule=0pt,
        top=20pt, bottom=5pt, left=25pt, right=5pt,
        % کلیدهای before و after حذف شدند تا تداخل ایجاد نکنند
        overlay={
            \fill[macRed] ([xshift=15pt,yshift=-12pt]frame.north west) circle (4pt);
            \fill[macYellow] ([xshift=27pt,yshift=-12pt]frame.north west) circle (4pt);
            \fill[macGreen] ([xshift=39pt,yshift=-12pt]frame.north west) circle (4pt);
            % نمایش نام فایل در سمت راست
            \ifstrempty{#1}{}{
                \node[text=white, anchor=east, font=\ttfamily\footnotesize] at ([xshift=-15pt,yshift=-12pt]frame.north east) {#1};
            }
        }
    }
}

% ------------------------------------------------------------------------------
% محیط برای کدهای درون‌خطی
% استفاده: \begin{code}[python]{نام فایل نمایشی} ... \end{code}
% ------------------------------------------------------------------------------
\newtcblisting{code}[2][python]{
    macbox={#2},
    listing only,
    minted language=#1,
    minted style=monokai, % استفاده از تم تیره
    minted options={
        fontsize=\footnotesize, 
        linenos, 
        numbersep=10pt, 
        breaklines=true,
        tabsize=4,
        baselinestretch=1.2
    }
}

% پیچیدن خودکار و امن محیط code در محیط latin با کمک etoolbox
\BeforeBeginEnvironment{code}{\begin{latin}}
\AfterEndEnvironment{code}{\end{latin}}

% ------------------------------------------------------------------------------
% دستور برای خواندن مستقیم از فایل (از پوشه codes)
% استفاده: \inputcode[python]{ch01.py}
% ------------------------------------------------------------------------------
% ابتدا دستور tcb را با یک نام داخلی تعریف می‌کنیم
\newtcbinputlisting{\myinputcode}[2][python]{
    macbox={#2},
    listing only,
    listing file={codes/#2},
    minted language=#1,
    minted style=monokai, % استفاده از تم تیره
    minted options={
        fontsize=\footnotesize, 
        linenos, 
        numbersep=10pt, 
        breaklines=true,
        tabsize=4,
        baselinestretch=1.2
    }
}

% سپس یک دستور پوششی (Wrapper) می‌سازیم که مشکل LTR را حل کند
\newcommand{\inputcode}[2][python]{%
    \begin{latin}%
    \myinputcode[#1]{#2}%
    \end{latin}%
}